\documentclass[11pt,oneside,reqno]{article}

\usepackage[T1]{fontenc}
\usepackage{amsmath,amssymb,amsthm}
\usepackage{mathpazo}
\usepackage[margin=1in]{geometry}
\usepackage{microtype}
\usepackage{tikz-cd}
\usetikzlibrary{arrows.meta,positioning}
\usepackage[hidelinks]{hyperref}
\hypersetup{
  pdftitle={Hochschild Homology in the Boutet de Monvel Calculus},
  pdfauthor={Changwei Zhou}
}

\DeclareMathOperator{\HH}{HH}
\DeclareMathOperator{\gr}{gr}
\DeclareMathOperator{\Res}{Res}
\DeclareMathOperator{\tr}{tr}
\newcommand{\C}{\mathbb C}
\newcommand{\Z}{\mathbb Z}
\newcommand{\cB}{\mathcal B}
\newcommand{\cJ}{\mathcal J}
\newcommand{\Bcs}{\mathfrak B_{\mathrm{cs}}}
\newcommand{\Splus}{\mathcal S_+}
\allowdisplaybreaks

\theoremstyle{plain}
\newtheorem{theorem}{Theorem}[section]
\newtheorem{proposition}[theorem]{Proposition}
\newtheorem{lemma}[theorem]{Lemma}
\newtheorem{corollary}[theorem]{Corollary}

\theoremstyle{definition}
\newtheorem{definition}[theorem]{Definition}
\newtheorem{example}[theorem]{Example}

\theoremstyle{remark}
\newtheorem{remark}[theorem]{Remark}
\newtheorem*{discussion}{Discussion}

\title{Hochschild Homology in the\\
Boutet de Monvel Calculus}
\author{Changwei Zhou}
\date{}

\begin{document}
\hypersetup{pageanchor=false}
\maketitle
\thispagestyle{empty}
\clearpage
\hypersetup{pageanchor=true}
\pagenumbering{roman}
\tableofcontents
\clearpage
\listoffigures
\clearpage
\pagenumbering{arabic}
\setcounter{page}{1}

\begin{abstract}
Let \(X\) be a compact connected \(n\)-manifold with boundary, \(n>1\).
We study the integer-order, type-zero Boutet de Monvel algebra modulo its
residual ideal, with a specified topologically filtered Hochschild complex.
The order spectral sequence has second page
\[
 E^2_k\cong
 H^{2n-k}(S^*X\times S^1_\sigma,
 S^*X|_{\partial X}\times S^1_\sigma;\C).
\]
The boundary term comes from the smooth Wiener--Hopf ideal and is essential:
discarding it gives the wrong trace space. A strong convergence and
\(E^2\)-collapse theorem would turn the displayed formula into the full
continuous Hochschild homology. We state that implication precisely, but do
not attribute it to the groupoid theorem of Benameur--Nistor, which does not
cover the full Boutet de Monvel matrix calculus. Independently of collapse,
the theorem of Fedosov--Golse--Leichtnam--Schrohe computes the Hausdorff
degree-zero group: it is one-dimensional and is detected by the boundary
noncommutative residue. Thus the manuscript separates the calculated page
from the one remaining global homological input.
\end{abstract}

\section{Scope and main statements}

The expression \emph{Hochschild homology of the Boutet de Monvel calculus}
does not specify an algebra. The answer changes if one keeps smoothing
operators, uses only order zero, admits arbitrary type, or changes the
locally convex completion. We use the classical integer-order algebra of
type zero, normalize its four entries with fixed order-reduction operators,
and divide by its full residual
ideal. This is the boundary analogue of the complete-symbol algebra on a
closed manifold.

Set
\[
 Y=S^*X=(T^*X\setminus0)/\mathbb R_+,
 \qquad Z=Y|_{\partial X}.
\]
Then \(Y\) is a compact manifold with boundary \(Z\). The circle
\(S^1_\sigma\) records integral symbol order; it is not an extra geometric
circle in \(T^*X\).

\begin{theorem}[The computed page]\label{thm:E2}
For the topologically filtered complete-symbol algebra defined in
Section~\ref{sec:algebra}, the order spectral sequence has
\begin{equation}\label{eq:E2-main}
 \boxed{
 E^2_k\cong
 H^{2n-k}(Y\times S^1_\sigma,Z\times S^1_\sigma;\C).}
\end{equation}
Equivalently,
\[
 E^2_k\cong H^{2n-k}(Y,Z;\C)
 \oplus H^{2n-k-1}(Y,Z;\C).
\]
\end{theorem}

\begin{theorem}[Conditional full calculation]\label{thm:conditional}
Assume that the order spectral sequence is strongly convergent and collapses:
\begin{equation}\label{eq:collapse}
 E^2(\Bcs)=E^\infty(\Bcs)
 \quad\text{and}\quad
 E^\infty(\Bcs)\Longrightarrow\HH_*^{\mathrm{cont}}(\Bcs).
\end{equation}
Then the associated graded of the induced homology filtration is canonically
given by the right side of \eqref{eq:E2-main}. After choosing linear
splittings,
\[
 \HH_k^{\mathrm{cont}}(\Bcs)
 \cong
 H^{2n-k}(Y\times S^1_\sigma,Z\times S^1_\sigma;\C)
\]
noncanonically as vector spaces. In particular, it vanishes for \(k>2n\).
\end{theorem}

\begin{theorem}[Unconditional degree zero]\label{thm:HH0}
For connected \(X\), \(n>1\), the Hausdorffized continuous degree-zero
homology is
\[
 \overline{\HH}^{\mathrm{cont}}_0(\Bcs)
 :=\Bcs/\overline{[\Bcs,\Bcs]}
 \cong\C.
\]
It is detected by the Fedosov--Golse--Leichtnam--Schrohe residue. This
statement does not require \eqref{eq:collapse}. Without a separate
closed-range theorem, the closure on the commutator space must be retained.
\end{theorem}

\begin{figure}[htbp]
\centering
\begin{tikzcd}[column sep=large,row sep=large]
 \Bcs \arrow[r,"\text{order filtration}"]
 & E^1 \arrow[r,"\text{Poisson boundary}"]
 & E^2 \arrow[d,"\text{convergence + collapse?}"]\\
 \HH_*^{\mathrm{cont}}(\Bcs)
 &H^{2n-*}(Y\!\times\!S^1_\sigma,Z\!\times\!S^1_\sigma)
  \arrow[l,dashed,"\text{assembly}"']
 &E^\infty\arrow[l,equals]
\end{tikzcd}
\caption[Proved and conditional parts of the calculation]{The symbol and
Poisson steps compute \(E^2\). Passing from \(E^2\) to the full homology
requires the separately stated convergence-and-collapse input.}
\end{figure}

\section{The algebra and topology}\label{sec:algebra}

Let \(E\to X\) and \(F\to\partial X\) be complex vector bundles. An operator
in Boutet de Monvel's calculus has the matrix form
\begin{equation}\label{eq:BdM-matrix}
 A=\begin{pmatrix}
       P_++G & K\\
       T     & S
    \end{pmatrix}:
 C^\infty(X;E)\oplus C^\infty(\partial X;F)
 \longrightarrow
 C^\infty(X;E)\oplus C^\infty(\partial X;F).
\end{equation}
Here \(P\) is a classical pseudodifferential operator on the double of
\(X\), its symbol satisfies the transmission condition, and
\(P_+=r^+Pe^+\). The other entries are, respectively, a singular Green,
potential, trace, and boundary pseudodifferential operator. We use type zero,
which is closed under composition and contains parametrices after the usual
order reductions \cite{Boutet,Grubb}.

The four blocks have shifted natural orders. Choose elliptic scalar order
reductions on \(X\) and \(\partial X\). Conjugating by their diagonal matrix
puts all blocks into one integer filtration. Different choices give
filtration-preserving Morita equivalences. Denote the resulting algebra by
\[
 \cB^{\mathbb Z}_0(X;E,F)=\bigcup_{m\in\Z}\cB^m_0(X;E,F).
\]
Its residual ideal \(\cB^{-\infty}\) consists of matrices whose entries have
smooth kernels in the appropriate interior, boundary, trace, and potential
spaces. Define
\begin{equation}\label{eq:Bcs}
 \Bcs(X;E,F)=
 \cB^{\mathbb Z}_0(X;E,F)/\cB^{-\infty}(X;E,F).
\end{equation}
Keeping \(\cB^{-\infty}\) produces a different extension problem and an
additional ordinary operator trace.

We use the topologically filtered chain space of Benameur--Nistor
\cite[Definition~1]{BenameurNistorFamilies}. Write
\(F^N_{N,N}\Bcs\) for their diagonal cofinal family: one first completes a
finite symbol-order window in the decreasing order direction and only then
lets the upper-order bound tend to infinity. In Hochschild degree \(q\), set
\begin{equation}\label{eq:chains}
 C_q^{\mathrm{tf}}(\Bcs)
 =\varinjlim_N
 \bigl(F^N_{N,N}\Bcs\bigr)^{\widehat\otimes_\pi(q+1)}.
\end{equation}
This notation packages the projective limits defining each completed
filtration piece; it is not the tensor power of one order quotient.
Multiplication is continuous on fixed pieces and sends a diagonal piece into
a later diagonal piece, so the Hochschild boundary extends to
\eqref{eq:chains}. We denote its homology by
\(\HH_*^{\mathrm{cont}}(\Bcs)\).

\begin{discussion}
The topology in \eqref{eq:chains} is part of the problem. Taking one
projective tensor power of the entire LF union need not make symbol
composition jointly continuous. Formula \eqref{eq:chains} keeps each chain at
finite upper order while allowing a complete classical asymptotic tail.
\end{discussion}

\begin{figure}[htbp]
\centering
\begin{tikzcd}[row sep=large,column sep=large]
 C^\infty(X;E)\arrow[r,"P_++G"]
 &C^\infty(X;E)
 &C^\infty(\partial X;F)\arrow[r,"K"]
 &C^\infty(X;E)\\
 C^\infty(X;E)\arrow[r,"T"]
 &C^\infty(\partial X;F)
 &C^\infty(\partial X;F)\arrow[r,"S"]
 &C^\infty(\partial X;F)
\end{tikzcd}
\caption[Interior and boundary blocks]{The diagonal blocks act in the
interior and on the boundary; the trace and potential blocks couple the two
spaces. The singular Green correction belongs to the interior block.}
\end{figure}

\section{The boundary symbol and the Wiener--Hopf ideal}

There are two compatible full symbol systems. The interior symbol is the
transmission symbol of \(P\) on \(T^*X\setminus0\). The boundary symbol is an
operator-valued classical symbol on \(T^*\partial X\setminus0\), acting on
\[
 \Splus\otimes E|_{\partial X}\oplus F,
 \qquad \Splus=\mathcal S(\overline{\mathbb R}_+).
\]
In product coordinates its leading term is
\begin{equation}\label{eq:boundary-symbol}
 a_\partial(x',\xi')=
 \begin{pmatrix}
 p_+(x',0,\xi',D_n)+g(x',\xi')&k(x',\xi')\\
 t(x',\xi')&s(x',\xi')
 \end{pmatrix}.
\end{equation}
The Toeplitz part is determined by the boundary restriction of the interior
symbol; \(g,k,t,s\) are independent boundary data.

Let \(\mathcal K^\infty(\Splus)\) be the algebra of Schwartz-kernel compact
operators. The smooth Wiener--Hopf fiber algebra has the linking ideal
\begin{equation}\label{eq:linking}
 \cJ^\infty=
 \begin{pmatrix}
  \mathcal K^\infty(\Splus)&\Splus\\
  \Splus'&\C
 \end{pmatrix}.
\end{equation}

\begin{lemma}[Smooth Morita reduction]\label{lem:morita}
The algebra \(\cJ^\infty\) is nuclear and \(H\)-unital, and
\[
 \HH_0^{\mathrm{cont}}(\cJ^\infty)\cong\C,
 \qquad
 \HH_q^{\mathrm{cont}}(\cJ^\infty)=0\quad(q>0).
\]
The degree-zero map is the full linking trace
\[
 \begin{pmatrix}K&u\\v&c\end{pmatrix}\longmapsto
 \operatorname{Tr}(K)+c.
\]
\end{lemma}

\begin{proof}
Choose a rank-one projection \(e\) with Schwartz range. The corner
\(e\cJ^\infty e\) is \(\C e\), and \(e\) is full. Schwartz nuclearity makes
the completed finite-rank factorization strict, while the standard
matrix-unit homotopy contracts the reduced bar complex. Thus
\(\cJ^\infty\) is smoothly Morita equivalent to \(\C\). The same argument
supplies the \(H\)-unitality required by continuous Hochschild excision
\cite{MeyerExcision}.
\end{proof}

The quotient of the Wiener--Hopf algebra by \(\cJ^\infty\) remembers the
scalar normal symbol. The transmission condition identifies it with the
restriction of the interior homogeneous symbol at \(Z\). Therefore the
order-graded Boutet de Monvel symbol algebra is a smooth pullback, not a
direct sum.

\begin{figure}[htbp]
\centering
\begin{tikzcd}[column sep=large,row sep=large]
 \gr\Bcs\arrow[r]\arrow[d]
 &C^\infty_{\mathrm{hom}}(T^*X\setminus0)
   \arrow[d,"i^*"]\\
 C^\infty_{\mathrm{hom}}
 (T^*\partial X\setminus0;\mathcal W^\infty)
 \arrow[r,"\text{normal quotient}"']
 &C^\infty_{\mathrm{hom}}
 (T^*X|_{\partial X}\setminus0)
\end{tikzcd}
\caption[Pullback structure of the graded boundary symbol]{The interior and
boundary symbols meet through the normal quotient. Hochschild excision turns
this pullback into a relative, rather than direct-sum, form complex.}
\end{figure}

\section{The first page}

Filter a chain by the sum of the upper orders of its factors. Classical
composition preserves the filtration. The leading product is the pointwise
pullback product above. This constructs the order exact couple and its
spectral sequence.

Use the mapping-cone model
\begin{equation}\label{eq:relative-cone}
 \Omega^j(Y,Z)=\Omega^j(Y)\oplus\Omega^{j-1}(Z),
 \qquad
 D(\alpha,\beta)=(d\alpha,i^*\alpha-d\beta).
\end{equation}
Its cohomology is \(H^j(Y,Z;\C)\). The second component is a boundary
transgression, not an independent boundary class.

\begin{proposition}[Boundary HKR identification]\label{prop:HKR}
Continuous Hochschild excision for \eqref{eq:linking}, followed by the HKR
map on the scalar quotient, identifies the first page with the homogeneous
mapping-cone complex for
\[
 (T^*X\setminus0,T^*X|_{\partial X}\setminus0).
\]
This identification is independent of collars and order reductions.
\end{proposition}

\begin{proof}
Work first in a product collar and trivialize \(E\) and \(F\). The pullback
of symbol algebras gives a short exact diagram whose boundary kernel is the
Schwartz linking ideal. Lemma~\ref{lem:morita} and continuous excision replace
that ideal by one scalar boundary complex. HKR identifies the scalar
interior and normal quotients with homogeneous differential forms. The
connecting map is restriction of the interior form to the boundary normal
symbol. Thus the resulting total complex is \eqref{eq:relative-cone}, before
radial reduction.

On overlaps, collar changes act by smooth Wiener--Hopf conjugations plus
lower-order terms. Inner conjugations are chain homotopic to the identity,
and lower-order terms do not alter this page. The local complexes are fine
sheaves over \(X\), so their local quasi-isomorphisms globalize. Order
reductions change the model by a full smooth Morita corner.
\end{proof}

\section{The Poisson differential}

For interior symbols \(a\) and \(b\),
\[
 \sigma_{m+m'-1}([A,B])=\frac1i\{a,b\}.
\]
Inserting this one-order-lower term in the Hochschild boundary and applying
HKR gives the Brylinski differential
\[
 \delta=[\iota_\pi,d].
\]
On the boundary, the regularized normal trace is not cyclic. Its commutator
defect is the restriction of the interior Poisson term. This produces the
second component \(i^*\alpha-d\beta\) in \eqref{eq:relative-cone}; it is also
the mechanism coupling the interior and boundary cochains in the boundary
index cocycle \cite{FGLS,BoltachevSavin}.

Let \(\omega\) be the canonical symplectic form on \(T^*X\setminus0\) and
\(\mu=\omega^n/n!\). Symplectic duality
\[
 *_\omega\alpha=\iota_{\pi^\sharp\alpha}\mu
\]
satisfies
\[
 *_\omega\delta=(-1)^{j+1}d*_\omega
 \quad\text{on \(j\)-forms}.
\]
Stokes' formula supplies the boundary transgression, so this identity extends
to the mapping cone. Hence the homology of the first differential is relative
de Rham cohomology with complementary degree.

Choose a positive fiber norm \(r\). For the Euler field
\(\mathcal E=r\partial_r\), Cartan's formula gives
\[
 d\iota_{\mathcal E}+\iota_{\mathcal E}d=\mathcal L_{\mathcal E}.
\]
Thus every nonzero weight is contracted by
\(w^{-1}\iota_{\mathcal E}\). A weight-zero relative form splits uniquely as
\[
 (\alpha_0,\beta_0)+\frac{dr}{r}\wedge(\alpha_1,\beta_1).
\]
Replacing \(dr/r\) by the positive generator of \(H^1(S^1_\sigma)\) yields
\[
 H^j_{\mathrm{hom}=0}\cong
 H^j(Y\times S^1_\sigma,Z\times S^1_\sigma).
\]
This proves Theorem~\ref{thm:E2}.

\section{The remaining convergence-and-collapse problem}

Benameur and Nistor compute complete symbols on differential groupoids and
prove convergence and collapse under explicit anchor and local-triviality
hypotheses \cite{BenameurNistorSymbols,BenameurNistorFamilies}. Their theorem
applies to the b-groupoid algebra in the companion manuscript. It does not
automatically apply to the matrix algebra \eqref{eq:BdM-matrix}.

The distinction is structural. Aastrup--Melo--Monthubert--Schrohe ask whether
the full Boutet de Monvel algebra is a groupoid pseudodifferential algebra and
construct a groupoid model for an important ideal, not an identification of
the whole calculus \cite{AMMS}. The boundary symbol contains Toeplitz, Green,
trace, and potential blocks; its pullback is not the scalar symbol algebra of
a Lie algebroid.

Completeness on each fixed upper-order piece constructs the spectral
sequence, but does not itself prove strong convergence after the strict
inductive limit. The unbounded order direction may carry a derived-limit
obstruction unless one proves a uniform regularity bound. Likewise, computing
the Poisson differential does not rule out higher differentials.

A proof of \eqref{eq:collapse} must:
\begin{enumerate}
 \item prove strict continuous excision uniformly in symbol order and a
       regular convergence bound;
 \item show that the Wiener--Hopf Morita homotopies preserve the complete
       order filtration;
 \item construct permanent cycles for all relative \(E^2\)-classes in a
       collar model;
 \item globalize those cycles without losing asymptotic completeness; and
 \item rule out all higher differentials.
\end{enumerate}
These tasks fit the Benameur--Nistor strategy, but none follows from Borel
summation alone. Once \eqref{eq:collapse} is proved, strong convergence
identifies \(E^\infty\) with the associated graded of homology and linear
splittings give Theorem~\ref{thm:conditional}.

\section{Degree zero and the boundary residue}

For a locally convex algebra,
\[
 \HH_0^{\mathrm{cont}}(A)
 =A/\operatorname{im}\bigl(b:C_1^{\mathrm{cont}}(A)\to A\bigr).
\]
Its Hausdorffization is \(A/\overline{[A,A]}\); its continuous dual is the
space of continuous traces.

For \(A\) as in \eqref{eq:BdM-matrix}, write \(p_{-n}\) for the degree
\(-n\) interior symbol, \((\tr_n g)_{1-n}\) for the degree \(1-n\) normal
trace of the Green symbol, and \(s_{1-n}\) for the corresponding boundary
symbol. With the standard cosphere densities, the
Fedosov--Golse--Leichtnam--Schrohe residue is
\begin{align}\label{eq:FGLS}
 \Res_{\mathrm{BdM}}(A)
 ={}&(2\pi)^{-n}\int_X\int_{S_x^*X}
       \tr_E p_{-n}(x,\xi)\,dS(\xi)\,dx\notag\\
 &+(2\pi)^{1-n}\int_{\partial X}\int_{S_{x'}^*\partial X}
 \left[
  \tr_E(\tr_n g)_{1-n}(x',\xi')
  +\tr_Fs_{1-n}(x',\xi')
 \right]dS(\xi')\,dx'.
\end{align}
It is coordinate invariant and vanishes on commutators
\cite{FGLS,Schrohe,GS}. The interior and boundary integrals are not
separately traces; their commutator defects cancel.

Under the connectedness and dimension hypotheses above, the FGLS trace
theorem says that every continuous trace on the complete-symbol algebra is a
multiple of \eqref{eq:FGLS}. The residue is nonzero. Since the continuous
dual separates points of a Hausdorff locally convex quotient,
\(\Bcs/\overline{[\Bcs,\Bcs]}\) is one-dimensional. This proves
Theorem~\ref{thm:HH0}.

The top relative group gives the same count:
\[
 H^{2n}(Y\times S^1_\sigma,Z\times S^1_\sigma)
 \cong H^{2n-1}(Y,Z)\cong\C.
\]
This is a consistency check, not a proof that the unclosed commutator range
is closed.

\section{Cyclic-cocycle check}

The boundary index cocycle independently confirms the mapping-cone geometry.
In the model of Boltachev--Savin, an interior cochain
\(\varphi_{\mathrm{int}}\) and a boundary regularized-trace cochain
\(\varphi_\partial\) satisfy
\[
 b\varphi_{\mathrm{int}}=0,\qquad
 b\varphi_\partial+cB\varphi_{\mathrm{int}}=0,\qquad
 B\varphi_\partial=0
\]
for the normalization constant \(c\) fixed by the normal Fourier convention
\cite{BoltachevSavin}. The middle identity is the cyclic counterpart of
\[
 D(\alpha,\beta)=(d\alpha,i^*\alpha-d\beta).
\]
Thus the boundary cochain transgresses the restricted interior class.
Treating it as an unrelated direct summand would incorrectly predict two
degree-zero residue traces.

\section{Checks and variants}

\subsection{Closed manifolds}

If \(\partial X=\varnothing\), then \(Z=\varnothing\) and
\[
 E^2_k\cong H^{2n-k}(S^*X\times S^1_\sigma;\C),
\]
the Brylinski--Getzler complete-symbol calculation. In that setting the
standard collapse theorem is available.

\subsection{Boundary components}

The pair \((Y,Z)\) includes every boundary component. For connected \(X\)
and \(n>1\), \(H^{2n-1}(Y,Z)\cong\C\), so the residue is one trace, not one
trace per boundary component. Lower relative groups retain the boundary
topology.

\subsection{Coefficient bundles}

Stabilization by \(E\) and \(F\) does not change the filtered calculation.
Locally this is matrix Morita invariance; globally the symbol modules form a
full equivalence bimodule. The matrix traces in \eqref{eq:FGLS} implement the
same Morita map.

\subsection{Unreduced and order-zero algebras}

The unreduced algebra belongs to the extension
\[
 0\longrightarrow\cB^{-\infty}
 \longrightarrow\cB^{\mathbb Z}_0
 \longrightarrow\Bcs\longrightarrow0.
\]
Its homology needs continuous excision and the homology of the chosen
residual ideal. The ordinary kernel trace lives in that ideal and is absent
from \(\Bcs\).

The \(C^*\)-closure of the order-zero calculus is the right setting for
operator \(K\)-theory, but it is not \eqref{eq:Bcs}. The work of
Melo--Nest--Schick--Schrohe and Melo--Schrohe--Schick describes its
\(K\)-theoretic symbol extension, not the Hochschild complex
\eqref{eq:chains} \cite{MNSS,MSS}.

\section{Conclusion}

The order calculation gives relative cohomology of the cosphere pair, with
the symbol-order circle, at \(E^2\). The boundary noncommutative residue
verifies the top-degree prediction and computes Hausdorff degree zero without
any collapse assumption. The sole remaining input for the all-degree formula
is the explicitly stated Boutet de Monvel convergence-and-collapse theorem.
This is narrower and more accurate than claiming that the existing groupoid
complete-symbol theorem already covers the boundary matrix calculus.

\begin{thebibliography}{99}

\bibitem{AMMS}
J.~Aastrup, S.~T. Melo, B.~Monthubert, and E.~Schrohe,
\emph{Boutet de Monvel's calculus and groupoids I},
J. Noncommut. Geom. \textbf{4} (2010), 313--329.

\bibitem{BenameurNistorSymbols}
M.-T. Benameur and V.~Nistor,
\emph{Homology of complete symbols and non-commutative geometry},
K-Theory \textbf{29} (2003), 201--223.

\bibitem{BenameurNistorFamilies}
M.-T. Benameur and V.~Nistor,
\emph{Homology of algebras of families of pseudodifferential operators},
J. Funct. Anal. \textbf{205} (2003), 1--36.

\bibitem{BoltachevSavin}
A.~V. Boltachev and A.~Yu. Savin,
\emph{Equivariant cyclic cocycles on the Boutet de Monvel symbol algebra},
arXiv:2201.09987.

\bibitem{Boutet}
L.~Boutet de Monvel,
\emph{Boundary problems for pseudo-differential operators},
Acta Math. \textbf{126} (1971), 11--51.

\bibitem{FGLS}
B.~V. Fedosov, F.~Golse, E.~Leichtnam, and E.~Schrohe,
\emph{The noncommutative residue for manifolds with boundary},
J. Funct. Anal. \textbf{142} (1996), 1--31.

\bibitem{Grubb}
G.~Grubb,
\emph{Functional Calculus of Pseudodifferential Boundary Problems},
second edition, Progress in Mathematics 65, Birkhauser, 1996.

\bibitem{GS}
G.~Grubb and E.~Schrohe,
\emph{Trace expansions and the noncommutative residue for manifolds with
boundary},
J. Reine Angew. Math. \textbf{536} (2001), 167--207.

\bibitem{MeyerExcision}
R.~Meyer,
\emph{Excision in Hochschild and cyclic homology without continuous linear
sections},
J. Homotopy Relat. Struct. \textbf{5} (2010), 269--303.

\bibitem{MNSS}
S.~T. Melo, R.~Nest, T.~Schick, and E.~Schrohe,
\emph{\(K\)-theory of Boutet de Monvel's algebra},
J. Reine Angew. Math. \textbf{561} (2003), 145--175.

\bibitem{MSS}
S.~T. Melo, E.~Schrohe, and T.~Schick,
\emph{\(C^*\)-algebra approach to the index theory of boundary value
problems},
J. Funct. Anal. \textbf{263} (2012), 3679--3714.

\bibitem{Schrohe}
E.~Schrohe,
\emph{Noncommutative residues, Dixmier's trace, and heat trace expansions on
manifolds with boundary},
in \emph{Geometric Aspects of Partial Differential Equations},
Contemp. Math. 242, American Mathematical Society, 1999, 161--186.

\end{thebibliography}

\end{document}
